\chapter{離散Fourier変換(DFT)}
    \section{基底}
        $d\in\naturalNumbers,N_l\in\naturalNumbers\;(l=1,2,\dots,d),\bm{k},\bm{n}\in\integers^d$とする。
        次式で定義される，$\bm{n}$に関する離散座標信号を$d$次元DFTの第$\bm{k}$基底ベクトルという。
        \[ W(\bm{k},\bm{n}) \coloneq \left(\prod_{l=1}^d N_l\right)^{-1/2} \exp i\left(\sum_{l=1}^d \frac{k_l n_l}{N_l}2\pi\right)\]

    \section{DFTの定義}
        \label{DFTの定義}
        $d\in\naturalNumbers,N_l\in\naturalNumbers\;(l=1,2,\dots,d),\bm{k}\in\integers^d$とする。
        $\Omega \coloneq \{0,1,\dots,N_1-1\}\times\{0,1,\dots,N_2-1\}\times\cdots\times\{0,1,\dots,N_d-1\}$とする。
        $f$を周期が$(N_1,N_2,\dots,N_d)$であるような離散座標信号$f: \integers^d\to\complexNumbers;\;\bm{n} = [n_1,n_2,\dots,n_d]^\top \mapsto f(\bm{n})$とするとき，次式で定義される，$\bm{k}$に関する離散座標信号を$f$の離散Fourier変換(Discrete Fourier Transform, DFT)という。
        \[ \DFTwithArg{f}{\bm{k}} \coloneq \sum_{\bm{n}\in\Omega} \conj{W(\bm{k},\bm{n})} f(\bm{n}) \]

    \section{Hermiteな離散座標信号のDFTは実数である}
        \label{Hermiteな離散座標信号のDFTは実数である}
        \begin{shadebox}
            $d,N_l,\bm{k},\Omega,f$の定義は\ref{DFTの定義}と同じものとする。
            $f$にさらにHermite性: $\conj{f(\bm{n})} = f(-\bm{n})$を要請するとき，$\DFTwithArg{f}{\bm{k}}$は実数となる。
        \end{shadebox}
        \begin{proof}
            \begin{align*}
                2\Im{\DFTwithArg{f}{\bm{k}}} &= \DFTwithArg{f}{\bm{k}} - \conj{\DFTwithArg{f}{\bm{k}}} \\
                &= \sum_{\bm{n}\in\Omega} \conj{W(\bm{k},\bm{n})}f(\bm{n}) - \sum_{\bm{n} \in \Omega}W(\bm{k},\bm{n})\conj{f(\bm{n})}
            \end{align*}
            ここで$\bm{n}_\text{M} \coloneq [N_1,\dots,N_d]^\top$とすると，
            \[ \conj{W(\bm{k},\bm{n})} = W(\bm{k},-\bm{n}) = W(\bm{k},\bm{n}_\text{M}-\bm{n}) \]
            また，$f$のHermite性の仮定より
            \[ \conj{f(\bm{n})} = f(-\bm{n}) = f(\bm{n}_\text{M}-\bm{n}) \]
            以上より
            \begin{align*}
                2\Im{\DFTwithArg{f}{\bm{k}}} &= \sum_{\bm{n}\in\Omega} W(\bm{k},\bm{n}_\text{M}-\bm{n})f(\bm{n}) - \sum_{\bm{n} \in \Omega}W(\bm{k},\bm{n})f(\bm{n}_\text{M}-\bm{n}) \\
                &= \sum_{\bm{n}\in\Omega} W(\bm{k},\bm{n}_\text{M}-\bm{n})f(\bm{n}) - \sum_{\bm{n}\in\Omega} W(\bm{k},\bm{n}_\text{M}-\bm{n})f(\bm{n}) \\
                &\phantom{=} (\{(\bm{n}, \bm{n}_\text{M}-\bm{n}) | \bm{n}\in\Omega\} = \{(\bm{n}_\text{M}-\bm{n}), \bm{n} | \bm{n}\in\Omega\}\text{を用いた}) \\
                &= 0
            \end{align*}
        \end{proof}

    \subsection{系: Hermiteな離散座標信号の IDFT は実数である}
        \ref{Hermiteな離散座標信号のDFTは実数である}と同様にして示せる。

    \section{巡回畳み込みのDFTはDFTの積に比例する}
        \begin{shadebox}
            $d,N_l,\bm{k},\Omega$の定義は\ref{DFTの定義}と同じものとする。
            $f,g$を周期が$(N_1,N_2,\dots,N_d)$であるような離散座標信号$f,g: \integers^d\to\complexNumbers;\;\bm{n} = [n_1,n_2,\dots,n_d]^\top \mapsto f(\bm{n}),g(\bm{n})$とするとき，次が成り立つ。
            \[ \DFTwithArg{f\circledast g}{\bm{k}} = \left(\prod_{l=1}^d N_l\right)^{1/2}\DFTwithArg{f}{\bm{k}}\DFTwithArg{g}{\bm{k}} \]
        \end{shadebox}
        \begin{proof}
            \quad\par
            $\bm{N} \coloneq [N_1,\dots,N_d]^\top$とする。
            \begin{align*}
                \DFTwithArg{f\circledast g}{\bm{k}} &= \sum_{\bm{n}\in\Omega}\conj{W(\bm{k},\bm{n})} \left(f\circledast g\right)(\bm{n}) = \sum_{\bm{n}\in\Omega}\conj{W(\bm{k},\bm{n})} \sum_{\bm{m}\in\Omega} f(\bm{m})g((\bm{n}-\bm{m})\%\bm{N}) \\
                &= \sum_{\bm{m}\in\Omega} f(\bm{m}) \sum_{\bm{n}\in\Omega} \conj{W(\bm{k},\bm{n})}g((\bm{n}-\bm{m})\%\bm{N}) \\
                &= \sum_{\bm{m}\in\Omega} f(\bm{m}) \sum_{\bm{n}\in\Omega} \left(\prod_{l=1}^d N_l\right)^{1/2} \conj{W(\bm{k},\bm{m})} \conj{W(\bm{k},\bm{n}-\bm{m})} g((\bm{n}-\bm{m})\%\bm{N}) \\
                &= \left(\prod_{l=1}^d N_l\right)^{1/2} \sum_{\bm{m}\in\Omega} \conj{W(\bm{k},\bm{m})}f(\bm{m}) \sum_{\bm{n}\in\Omega} \conj{W(\bm{k},(\bm{n}-\bm{m})\%\bm{N})} g((\bm{n}-\bm{m})\%\bm{N}) \\
                &= \left(\prod_{l=1}^d N_l\right)^{1/2} \sum_{\bm{m}\in\Omega} \conj{W(\bm{k},\bm{m})}f(\bm{m}) \sum_{\bm{n}\in\Omega} \conj{W(\bm{k},\bm{n})} g(\bm{n}) \\
                &= \left(\prod_{l=1}^d N_l\right)^{1/2}\DFTwithArg{f}{\bm{k}}\DFTwithArg{g}{\bm{k}}
            \end{align*}
        \end{proof}

    \section{巡回相関のDFT}
    \label{巡回相関のDFT}
        \begin{shadebox}
            $d,N_l,\bm{k},\Omega$の定義は\ref{DFTの定義}と同じものとする。
            $f,g$を周期が$(N_1,N_2,\dots,N_d)$であるような離散座標信号$f,g: \integers^d\to\complexNumbers;\;\bm{n} = [n_1,n_2,\dots,n_d]^\top \mapsto f(\bm{n}),g(\bm{n})$とする。
            $f$と$g$の巡回相関を$R_{f,g}(n) = \cycCorrel{f}{g}(\bm{n})$とする。
            このとき，次が成り立つ。
            \[ \DFTwithArg{R_{f,g}}{\bm{k}} = \left(\prod_{l=1}^d N_l\right)^{1/2}\DFTwithArg{f}{\bm{k}}\conj{\DFTwithArg{g}{\bm{k}}} \]
        \end{shadebox}
        \begin{proof}
            \quad\par
            $\bm{N} \coloneq [N_1,\dots,N_d]^\top$とする。
            \begin{align*}
                \DFTwithArg{R_{f,g}}{\bm{k}} &= \sum_{\bm{n}\in\Omega} \conj{W(\bm{k},\bm{n})} \sum_{\bm{m}\in\Omega} f(\bm{m})\conj{g((\bm{m}-\bm{n})\%\bm{N})} \\
                &= \sum_{\bm{m}\in\Omega} f(\bm{m}) \sum_{\bm{n}\in\Omega} \conj{W(\bm{k},\bm{n})}\conj{g((\bm{m}-\bm{n})\%\bm{N})} \tag{1}
            \end{align*}
            ここで$\conj{W(\bm{k},\bm{n})}$を変形して次式を得る。
            \begin{align*}
                \conj{W(\bm{k},\bm{n})} &= W(\bm{k},-\bm{n}) = \left(\prod_{l=1}^d N_l\right)^{1/2}W(\bm{k},\bm{m}-\bm{n})W(\bm{k},-\bm{m}) \\
                &= \left(\prod_{l=1}^d N_l\right)^{1/2}W(\bm{k},(\bm{m}-\bm{n})\%\bm{N})\conj{W(\bm{k},\bm{m})}
            \end{align*}
            これを式(1)に適用して次式を得る。
            \begin{align*}
                \DFTwithArg{R_{f,g}}{\bm{k}} &= \left(\prod_{l=1}^d N_l\right)^{1/2} \sum_{\bm{m}\in\Omega} \conj{W(\bm{k},\bm{m})}f(\bm{m}) \sum_{\bm{n}\in\Omega}\conj{\conj{W(\bm{k},(\bm{m}-\bm{n})\%\bm{N})}g((\bm{m}-\bm{n})\%\bm{N})} \\
                &= \left(\prod_{l=1}^d N_l\right)^{1/2}\DFTwithArg{f}{\bm{k}}\conj{\DFTwithArg{g}{\bm{k}}}
            \end{align*}
        \end{proof}
    \section{巡回畳み込みのDFT}
        \begin{shadebox}
            $d,N_l,\bm{k},\Omega$の定義は\ref{DFTの定義}と同じものとする。
            $f,g$を周期が$(N_1,N_2,\dots,N_d)$であるような離散座標信号$f,g: \integers^d\to\complexNumbers;\;\bm{n} = [n_1,n_2,\dots,n_d]^\top \mapsto f(\bm{n}),g(\bm{n})$とするとき，次が成り立つ。
            \[ \DFTwithArg{f\circledast g}{\bm{k}} = \left(\prod_{l=1}^d N_l\right)^{1/2}\DFTwithArg{f}{\bm{k}}\DFTwithArg{g}{\bm{k}} \]
        \end{shadebox}
        \begin{proof}
            \quad\par
            \ref{巡回相関のDFT} と同じ要領で示せる。
        \end{proof}
    \section{GaussianノイズのDFT}
        \begin{shadebox}
            $F(n) \in \complexNumbers\;(n=0,1,\dots,N-1)$は互いに独立で，複素正規分布$N(0,\sigma^2)$に従うとする$\left(\text{p}(f) = \frac{1}{2\pi\sigma^2}\exp\frac{-\Re{f}^2-\Im{f}^2}{2\sigma^2} = \frac{1}{2\pi\sigma^2}\exp\frac{-|f|^2}{2\sigma^2}\right)$。
            これのDFTを$G(k) = \DFTwithArg{F}{k}$とするとき，$G(k)\;(k=0,1,\dots,N-1)$もまた互いに独立で，複素正規分布$N(0,\sigma^2)$に従う。
        \end{shadebox}
        \begin{proof}
            \[ P \in \complexNumbers^{N\times N}, P_{k,n} \coloneq W(k,n) \coloneq \frac{1}{\sqrt{N}}\exp i\frac{kn}{N}2\pi\;(k,n \in \{0,1,\dots,N-1\}) \]
            \[ \bm{F} \coloneq [F(0), F(1), \dots, F(N-1)]^\top,\bm{G} \coloneq [G(0), G(1), \dots, G(N-1)]^\top \]
            と定義すると
            \[ \bm{G} = P^*\bm{F} \]
            となる。
            \begin{align*}
                \Pr{\bm{G} = \bm{g} \in \complexNumbers^N} &= \Pr{P^*\bm{F} = \bm{g}} = \Pr{\bm{F} = P\bm{g}} \\
                &= \prod_{i=0}^{N-1}\frac{1}{2\pi\sigma^2} \exp\frac{-|(P\bm{g})[i]|^2}{2\sigma^2} = \left(\prod_{i=0}^{N-1}\frac{1}{2\pi\sigma^2}\right) \exp \sum_{i=0}^{N-1} \frac{-|(P\bm{g})[i]|^2}{2\sigma^2} \\
                &= \left(\prod_{i=0}^{N-1}\frac{1}{2\pi\sigma^2}\right) \exp \frac{-\norm{P\bm{g}}_2^2}{2\sigma^2} = \left(\prod_{i=0}^{N-1}\frac{1}{2\pi\sigma^2}\right) \exp \frac{-\norm{\bm{g}}_2^2}{2\sigma^2} \\
                &= \prod_{i=0}^{N-1}\frac{1}{2\pi\sigma^2} \exp \frac{-|g_i|^2}{2\sigma^2}
            \end{align*}
        \end{proof}
    \section{DTFT を DFT で表す}
        \newcommand*{\XprdDFT}{\mathring{X}_\text{DFT}}
        有限長の信号は周期拡張によって DFT と一対一対応し，かつ DTFT とも一対一対応するので，信号が有限長であれば DTFT と DFT は一対一対応する。
        一次元の場合について示すが，多次元の場合にも容易に拡張できる。
        \begin{shadebox}
            $L\in\naturalNumbers$ とする。
            台が $I\coloneq\{0,1,\dots,L-1\}$ に含まれる離散時間信号 $x$ を周期 $L$ で周期的に拡張した信号を $\mathring{x}$ とする。
            $\XDTFT,\XprdDFT$ をそれぞれ $x$ の DTFT, $\mathring{x}$ の DFT とする。
            次式が成り立つ。
            ここに $F$ は正規化周波数であり，$x$ のサンプル周期 $\Ts$，周波数 $f$ を用いて $F = f\Ts$ と定義される。
            また $F = l/L\;(l\in\integers)$ に於ける値は $F$ がその値に近付く極限の値として定義される。
            \[ \XDTFT(F) = \frac{1}{\sqrt{L}}\sum_{k=0}^{L-1}\bracks*{\exp i(L-1)(k/L-F)\pi}\frac{\sin L(k/L-F)\pi}{\sin(k/L-F)\pi}\XprdDFT(k) \]
            とくに $F = l/L\;(l\in\integers)$ については次式が成り立つ。
            \[ \XDTFT(l/L) = \sqrt{L}\XprdDFT(l) \]
        \end{shadebox}
        \begin{proof}
            \begin{align*}
                \XDTFT(F) &= \sum_{m=0}^{L-1}\IDFT{\XprdDFT}(m)\exp(-i 2\pi F m) \\
                &= \sum_{m=0}^{L-1}\frac{1}{\sqrt{L}}\sum_{k=0}^{L-1}\XprdDFT(k)\parens*{\exp i\frac{m}{L}k 2\pi}\exp(-i 2\pi F m) \\
                &= \frac{1}{\sqrt{L}}\sum_{k=0}^{L-1}\XprdDFT(k)\sum_{m=0}^{L-1}\exp i m(k/L-F)2\pi \tag{1}
            \end{align*}
            $F = l/L\;(l\in\integers)$ のとき $\sum_{m=0}^{L-1}\exp i m(k/L-F)2\pi = \delta_{k,l}$ となる（等比数列の和の公式から容易に導ける）ので，$\XDTFT(l/L) = \sqrt{L}\XprdDFT(l)$ である。
            \par
            $F \neq l/L\;(l\in\integers)$ のとき，次式が成り立つ。
            \begin{align*}
                &\phantom{=} \sum_{m=0}^{L-1}\exp i m(k/L-F)2\pi = \frac{1 - \exp i L(k/L-F)2\pi}{1 - \exp i(k/L-F)2\pi} \\
                &= \frac{\exp iL(k/L-F)\pi}{\exp i(k/L-F)\pi}\times\frac{\exp(-i L(k/L-F)\pi) - \exp i L(k/L-F)\pi}{\exp(-i(k/L-F)\pi) - \exp i(k/L-F)\pi} \\
                &= [\exp i(L-1)(k/L-F)\pi]\times\underbrace{\frac{\sin L(k/L-F)\pi}{\sin(k/L-F)\pi}}_{(2)}
            \end{align*}
            これを式 (1) に適用して主張が得られる。
        \end{proof}
        数値計算に於いて，$\abs{k/L-F}\pi\ll 1$ のときに式 (2) の分数を安定的に計算するには，Taylor 展開 $(\sin Lx)/\sin x \approx L(1-(L^2-1)x^2/6)$ を用いるとよい。
